%
% cauchyverteilung.tex -- 
%
% (c) 2021 Prof Dr Andreas Müller, OST Ostschweizer Fachhochschule
%
\documentclass[tikz]{standalone}
\usepackage{amsmath}
\usepackage{times}
\usepackage{txfonts}
\usepackage{pgfplots}
\usepackage{csvsimple}
\usetikzlibrary{arrows,intersections,math,calc}
\definecolor{darkred}{rgb}{0.8,0,0}
\begin{document}
\def\skala{1}
\def\R{4.5}
\def\r{0.9}
\def\w{120}
\begin{tikzpicture}[>=latex,thick,scale=\skala]

\node at (\R,\R) [above right] {$\mathbb{C}$};

% Koordinatenachsen
\draw[->] (0,-5.1) -- (0,5.3) coordinate[label={right:$\operatorname{Im}$}];
\fill[color=gray!40,opacity=0.7] (0,0) -- (0:\r) arc (0:\w:\r) -- cycle;
\node at (0,0) [above right] {$\varphi$};
\draw[->] (-5.1,0) -- (5.3,0) coordinate[label={$\operatorname{Re}$}];

% Radius R und Endpunkt
\draw[->,color=black] (0,0) -- (\w:\R);
\node at (\w:{0.5*\R}) [left] {$R$};
\node at ($(120:\R)+(0.4,0)$) [above left] {$R(\cos\varphi+i\sin\varphi)$};
\fill[color=darkred] (\w:\R) circle[radius=0.05];

% \gamma_2
\draw[color=darkred,line width=1.4pt] (-\R,0) -- (\R,0) arc(0:180:\R);
\draw[->,color=darkred,line width=1.4pt] (-\R,0.1) -- (-\R,0);
\node[color=darkred] at (45:\R) [above right] {$\gamma_2$};

% \bar\gamma_2
\draw[->,color=darkred!40,line width=1.4pt] (-\R,-0.1) -- (-\R,0);
\draw[color=darkred!40,line width=1.4pt] (-\R,0) arc(-180:0:\R);
\node[color=darkred!40] at (-45:\R) [below right] {$\bar{\gamma}_2$};

% \gamma_1
\draw[color=darkred,line width=1.4pt] (-\R,0) -- (\R,0);
\node[color=darkred] at ({0.5*\R},0) [above] {$\gamma_1$};

% Punkte +i und -i
\coordinate (plusI) at (0,1);
\coordinate (minusI) at (0,-1);

\fill[color=white] (plusI) circle[radius=0.08];
\draw[color=blue] (plusI) circle[radius=0.08];
\node[color=blue] at (plusI) [above right] {$+i$};

\fill[color=white] (minusI) circle[radius=0.08];
\draw[color=blue] (minusI) circle[radius=0.08];
\node[color=blue] at (minusI) [below right] {$-i$};

% Beschriftung x-Achse
\node at (-\R,0) [below left] {$-R\mathstrut$};
\node at (\R,0) [below right] {$R\mathstrut$};

% Verbindungspunkte der Kurventeile
\fill[color=darkred] (-\R,0) circle[radius=0.05];
\fill[color=darkred] (\R,0) circle[radius=0.05];

\end{tikzpicture}
\end{document}

